% Part: counterfactuals
% Chapter: minimal-change-semantics
% Section: introduction

\documentclass[../../../include/open-logic-section]{subfiles}

\begin{document}

\olfileid{cnt}{min}{int}

\olsection{ভূমিকা}

স্টালনেকার ও লুইস ``দেশলাই-কাঠিটি ঘষা হলে জ্বলে উঠত''-এর
মতো প্রতিবাস্তব শর্তবচনের ব্যাখ্যা প্রস্তাব করেন। এমন
বাক্যের সত্যতার শর্ত কীভাবে ঠিক বুঝতে হবে, সেটিই তাঁদের
প্রস্তাবের লক্ষ্য। উভয়ের মূল ধারণা: প্রতিবাস্তব শর্তবচনটি
সত্য কি না বিচার করতে সেই সম্ভাব্য জগৎগুলি বিবেচনা করতে
হবে যেগুলিতে পূর্বপক্ষ সত্য, অথচ বাস্তব জগৎ থেকে পার্থক্য
ন্যূনতম। সেই জগৎগুলিতে অনুপক্ষ সত্য হলে প্রতিবাস্তব
শর্তবচনটিও সত্য। ধরা যাক, হাতে একটি দেশলাই-কাঠি ও
দেশলাই-বাক্স নিয়ে দাঁড়িয়ে আছি। বাস্তব জগতে আমি শুধু
সেগুলির দিকে তাকিয়ে ভাবছি, কাঠিটি ঘষলে কী হতো।
বাস্তব জগৎ থেকে ন্যূনতম পরিবর্তনে কাঠিটি ঘষা হয় এমন
জগতে আমি কাজটি করার সিদ্ধান্ত নিই ও কাঠিটি ঘষি।
পরিবর্তনটি ন্যূনতম, কারণ অপ্রাসঙ্গিক অতিরিক্ত বদল ঘটে না:
% BN-SRC-767: পূর্বপক্ষের বদল ও তার প্রাসঙ্গিক ফলগুলি থাকে; আক্ষরিকভাবে অন্য সব তথ্য অপরিবর্তিত নয়।
কাঠিটি ঘষা এবং তার ফলে জ্বলে ওঠার মতো প্রাসঙ্গিক পরিবর্তনগুলি
অবশ্যই থাকে। তবে একই সঙ্গে
আমি লাফিয়ে উঠি না, কাঠিটি ঘষে আমার চুলে আগুন লাগে
না, আঙুলের সমস্ত শক্তি হঠাৎ হারাই না, জল-পিস্তল থেকে
ছোড়া জলে ভিজে যাই না ইত্যাদি। ওই বিকল্প পরিস্থিতিতে
কাঠিটি জ্বলে ওঠে। তাই কাঠিটি ঘষা হলে জ্বলে উঠত—এটি
সত্য।

এই স্বজ্ঞাগত ব্যাখ্যাকে প্রতিবাস্তব শর্তবচনের যুক্তিবিদ্যার
আনুষ্ঠানিক অর্থতত্ত্বের সঙ্গে যুক্ত করা যায়। লুইস এর
জন্য ``$\boxright$'' এবং স্টালনেকার ``$>$'' চিহ্ন
প্রবর্তন করেন। আমরা $\cif$ ব্যবহার করব এবং বচনমূলক
যুক্তিবিদ্যায় একে দ্বিপদী সংযোজক হিসেবে যোগ করব। ফলে
$!A \lif !B$ আকারের !!{formula}s-এর পাশাপাশি
$!A \cif !B$ আকারের !!{formula}s-ও থাকবে। মোডাল
যুক্তিবিদ্যার সম্পর্কমূলক অর্থতত্ত্বের মতোই আনুষ্ঠানিক
অর্থতত্ত্বে জগৎ-ভিত্তিক মডেলে !!{formula}s মূল্যায়িত হয়।
$\mSat{M}{!A \cif !B}[w]$-এর পরিতৃপ্তি-শর্ত
কিছু (অন্য) জগৎ~$w'$-তে $\mSat{M}{!A}[w']$ এবং
$\mSat{M}{!B}[w']$ দিয়ে নির্ধারিত। কোন~$w'$?
স্বজ্ঞাগতভাবে, যেখানে~$\mSat{M}{!A}[w']$ সত্য, সেগুলির মধ্যে~$w$-র
সবচেয়ে কাছের জগৎ বা জগৎগুলি। তাই মডেলে ``নিকটত্ব''
সম্পর্কের তথ্যও রাখতে হবে।

প্রতিবাস্তব পরিস্থিতি ছবিতে দেখানোর একটি শিক্ষণীয়
পদ্ধতি লুইস দেন। প্রতিটি সম্ভাব্য জগৎকে কেন্দ্র করে
অন্য জগৎগুলিকে ধারণকারী পরস্পর-অন্তর্ভুক্ত গোলকের
সমষ্টি থাকে; ছবিতে এগুলিকে সমকেন্দ্রিক বৃত্ত দিয়ে
দেখাই। দুটি গোলকের মাঝখানের জগৎগুলি কেন্দ্রের জগৎ
থেকে একই দূরত্বে; ভিতরের গোলকে থাকা জগৎগুলি আরও
কাছে এবং বাইরের গোলকে থাকা জগৎগুলি আরও দূরে।
\begin{center}
\begin{tikzpicture}[scale=.7]\small
  \spheresystem{5}
  \draw (0,0) node {$w$};
  \propositionintersect{0}{4}{40}{3.5}
  {\tikzset{proposition/.append={smooth,tension=1.4}}
  \proposition[shift={(1.6,-1)}]{90}{1}{30}{4}}
  \path (1.5,3) node[below] {$\formula{B}$};
  \path (3,0) node {$\formula{A}$};
\end{tikzpicture}
\end{center}
% BN-SRC-768: ক্ষুদ্রতম বলতে পূর্বপক্ষ-স্বীকারক গোলক; এমন গোলক থাকা এই ছবির সীমা।
এই ছবির মতো ক্ষুদ্রতম পূর্বপক্ষ-স্বীকারক গোলক থাকলে
নিকটতম $!A$-জগৎ বলতে সেই সব জগৎ~$w'$-কে বোঝায়, যেখানে~$!A$ সত্য
এবং যেগুলি কেন্দ্রীয় জগৎ~$w$-কে ঘিরে সবচেয়ে ছোট পূর্বপক্ষ-স্বীকারক গোলকের মধ্যে আছে
(ছবির ধূসর অংশ)। স্বজ্ঞাগতভাবে, সব নিকটতম
$!A$-জগতে $!B$ সত্য হলে $w$-তে $!A \cif !B$
পরিতৃপ্ত। এরূপ ক্ষুদ্রতম গোলক না-থাকলেও পরের বিভাগের
গোলক-সংজ্ঞায় প্রতিবাস্তব শর্তবচনের সত্যতা নির্ধারিত হয়।

\end{document}
